\documentclass[11pt,a4paper]{article}

% Packages
\usepackage[utf8]{inputenc}
\usepackage[T1]{fontenc}
\usepackage[margin=0.75in]{geometry}
\usepackage{enumitem}
\usepackage{hyperref}
\usepackage{titlesec}
\usepackage{xcolor}

% Formatting
\pagestyle{empty}
\setlength{\parindent}{0pt}
\setlist[itemize]{leftmargin=*, noitemsep, topsep=3pt}

% Section formatting
\titleformat{\section}{\Large\bfseries}{}{0em}{}[\titlerule]
\titlespacing*{\section}{0pt}{12pt}{6pt}

\titleformat{\subsection}[runin]{\bfseries}{}{0em}{}
\titlespacing*{\subsection}{0pt}{8pt}{0pt}

% Colors
\definecolor{linkcolor}{RGB}{0,102,204}
\hypersetup{
    colorlinks=true,
    urlcolor=linkcolor
}

\begin{document}

% Header
\begin{center}
    {\LARGE \textbf{ERIC MAYEUX}}\\
    \vspace{4pt}
    {\large Chargé de projet sénior TI | Senior IT Project Manager}\\
    \vspace{4pt}
    Montréal, QC, Canada\\
    \href{mailto:mayeux.e@gmail.com}{mayeux.e@gmail.com} | 438-884-7645
\end{center}

\vspace{8pt}

\section*{Profil Professionnel}
\noindent
\textbf{Chargé de projet sénior TI} avec 10+ ans d'expérience en \textbf{gestion de projets TI} dans le secteur bancaire et financier. Expertise confirmée en \textbf{pilotage d'initiatives stratégiques}, gouvernance de projets complexes et \textbf{alignement stratégique} des parties prenantes. Maîtrise des méthodologies \textbf{Agile (Scrum, SAFe) et Waterfall} — approche hybride selon contexte. \textbf{Certifié PMP} (en cours), CSM et AWS Cloud Practitioner. Bilinguisme français/anglais.

\section*{Compétences Techniques}

\textbf{Gestion de projets TI:} Pilotage d'initiatives stratégiques, gestion de programmes multi-projets, gouvernance, alignement stratégique, amélioration continue, gestion des risques et enjeux

\textbf{Leadership \& Influence:} Mobilisation des parties prenantes, négociation, gestion des conflits, coaching d'équipes multidisciplinaires, Mentorat, Recrutement

\textbf{Méthodologies:} Agile, Scrum, SAFe, Kanban, Waterfall, approche hybride, PMP, PMI-ACP

\textbf{Outils de planification \& suivi:} Jira, MS Project, Confluence, Datadog, Splunk, tableaux de bord KPI

\textbf{CI/CD \& DevOps:} Jenkins, Docker, Kubernetes, GitHub Actions, GitLab CI/CD, Harness, ArgoCD

\textbf{Intelligence Artificielle:} Agents IA, Playwright MCP, Crew AI, RPA, automatisation

\textbf{Architecture \& Cloud:} Certifié AWS Cloud Practitioner, Architectures Distribuées, Big Data

\vspace{6pt}

\section*{Réalisations}
\begin{itemize}
    \item \textbf{Pilotage d'initiatives stratégiques} à la BNC : coordination de releases de A à Z (estimation, planification, coordination, mise en production)
    \item Mise en place de \textbf{gouvernances pérennes} multi-équipes avec alignement stratégique des parties prenantes
    \item Gestion de programmes regroupant plusieurs projets sous une \textbf{gouvernance commune} (contexte SAFe Release Train)
    \item Forte capacité d'analyse et d'\textbf{amélioration continue} : audit qualité, état des lieux, optimisation des pratiques
    \item Mise en place d'agents IA pour l'\textbf{automatisation de tâches de gestion}, d'analyse et d'administratif
    \item Encadrement de 5+ professionnels TI dans un contexte agile SAFe international (Thaïlande)
\end{itemize}

\section*{Expérience Professionnelle}

\subsection*{Scrum Master -- Gestion de projet | Project Manager} \hfill \textit{Novembre 2025 -- Présent}\\
\textbf{Banque Nationale du Canada (BNC)} -- Montréal, QC\\
Pilote des projets TI stratégiques dans un contexte bancaire en \textbf{gestion de projets TI} et transformation agile.
\begin{itemize}
    \item Encadre et coordonne des équipes IT multidisciplinaires (QA, Dev, PO, Analystes) vers les \textbf{objectifs stratégiques}
    \item \textbf{Pilote les releases} : estimation, planification, coordination, arrimage inter-équipes, mise en production
    \item Identifie et gère les \textbf{risques et enjeux}, propose des plans de mitigation, adapte la communication aux parties prenantes
    \item Crée des tableaux de bord pour le suivi des indicateurs de performance et la \textbf{gouvernance} des projets
    \item Volet innovation : développe des agents IA pour l'\textbf{automatisation} de tâches de gestion et d'administratif
    \item \textbf{Coaching agile} et amélioration continue : formations, ateliers, adaptation des pratiques au contexte
\end{itemize}

\subsection*{Intégrateur Sénior Qualité | Senior Quality Integrator} \hfill \textit{Octobre 2023 -- Novembre 2025}\\
\textbf{Banque Nationale du Canada (BNC)} -- Montréal, QC\\
Pilotage de la pratique qualité avec \textbf{leadership d'équipes techniques} et \textbf{gestion de projets TI} en contexte SAFe.
\begin{itemize}
    \item Encadre, accompagne, mentor les QA et Intégrateurs ; recrutement et formation continue
    \item Conçoit et exécute des \textbf{stratégies pan-train} dans un contexte \textbf{multi-sites} (Thaïlande)
    \item Met en place des métriques de qualité et tableaux de bord via Datadog, Splunk, Jira
    \item Conduit des \textbf{POC} et projets d'innovation : Playwright MCP, Crew AI, hackathon de test
    \item CI/CD/CT : met en place des pipelines intégrant l'exécution continue des tests
\end{itemize}

\subsection*{Intégrateur Qualité -- SDET | Quality Integrator -- Software Development Engineer in Test} \hfill \textit{Janvier 2022 -- Octobre 2023}\\
\textbf{Banque Nationale du Canada (BNC)} via Groupe Astek -- Montréal, QC\\
Intégrateur qualité en environnement agile — Gestion de Patrimoine \& ECM.
\begin{itemize}
    \item Conception, exécution et suivi de la \textbf{stratégie de test} et gestion des risques qualité
    \item Mise en place des métriques : pyramide de test, couverture, analyse statique de sécurité
    \item Stack : Selenium, AppliTools, RestAssured, JMeter, Gatling, K6
\end{itemize}

\subsection*{QA Automaticien | QA Automation Engineer} \hfill \textit{Juin 2021 -- Décembre 2021}\\
\textbf{AODocs} -- Paris, France
\begin{itemize}
    \item Automatisation des tests : Selenium, Java 8, JUnit 5, Page Object Model, Karate (API)
    \item POC tests mobiles Flutter ; rapports Allure
\end{itemize}

\subsection*{Automaticien CI/CD | DevOps Automation Engineer} \hfill \textit{Janvier 2020 -- Juin 2021}\\
\textbf{ENGIE DIGITAL} via Groupe Hénix -- Paris, France
\begin{itemize}
    \item Création de la chaîne CI/CD : Jenkins, Docker, GitHub
    \item \textbf{Administration Jira}, pilotage projets et mise en place de KPIs
\end{itemize}

\subsection*{Product Owner} \hfill \textit{Juin 2019 -- Septembre 2019}\\
\textbf{Hénix} -- Montrouge, France
\begin{itemize}
    \item Gestion du backlog, recette de plugins Jira/Squash, support clients SaaS
\end{itemize}

\subsection*{Développeur Logiciel | Software Developer} \hfill \textit{Mai 2017 -- Mai 2019}\\
\textbf{MEDIAPOST} via Groupe Hénix -- Paris, France
\begin{itemize}
    \item Développement d'applications back office, webservices SOAP/REST, scripts PL/SQL, Sybase
\end{itemize}

\subsection*{Consultant QA \& Automatisation | QA Automation Consultant} \hfill \textit{Avril 2016 -- Avril 2017}\\
\textbf{ENGIE IT} via Groupe Hénix -- Saint-Ouen, France
\begin{itemize}
    \item Pilotage opérationnel des environnements, administration HP ALM/UFT, gestion des incidents
\end{itemize}

\section*{Certifications \& Formations}

\textbf{PMP Certification} (en cours) \hfill \textit{2026}

Project Management Professional -- Project Management Institute (PMI)

\textbf{AWS Cloud Practitioner} \hfill \textit{2026}

Amazon Web Services (AWS)

\textbf{Certified Scrum Master (CSM)} \hfill \textit{2024}

Scrum Alliance

\textbf{Jira ACP-600 -- Project Administration in Jira Server} \hfill \textit{2019}

Atlassian Certified Professional

\textbf{Cursus Automaticien et DevOps} \hfill \textit{2019}\\
\textit{AFCEPF -- Montrouge, France}
\begin{itemize}
    \item Selenium, Jenkins, Docker, Kubernetes, Ansible, SoapUI, NeoLoad
    \item \textbf{Certifications:} NeoLoad Professional Certified, Certified Selenium Professional
\end{itemize}

\textbf{Cursus Architecture logiciel \& Analyste informaticien - Certifié Master II} \hfill \textit{2017}\\
\textit{AFCEPF-HENIX -- Malakoff, France}

\section*{Formation Académique}

\textbf{Licence en Gestion (Bachelor's Degree in Management)} \hfill \textit{2011}\\
École Supérieure de Gestion (E.S.G) -- Paris, France

\section*{Compétences Linguistiques}

\textbf{Français:} Langue maternelle (Native)\\
\textbf{Anglais:} Niveau Avancé (Advanced / Professional Working Proficiency)

\end{document}
